\section{Diffusion Probabilistic Models}
\label{sec:diffusion_theory}

Diffusion models are likelihood-based generative models that learn to synthesize data by reversing a gradual noising process. Unlike GANs, diffusion models optimize a tractable variational objective and provide stable training with high sample quality.

\subsection{Forward Diffusion Process}

Let $\mathbf{x}_0 \sim p_{\text{data}}(\mathbf{x})$ denote a real image.

The forward process gradually adds Gaussian noise over $T$ timesteps:

\begin{equation}
q(\mathbf{x}_t \mid \mathbf{x}_{t-1}) =
\mathcal{N}\left(
\sqrt{1 - \beta_t}\mathbf{x}_{t-1},
\beta_t I
\right)
\end{equation}

where:
\begin{itemize}
    \item $\beta_t \in (0,1)$ is a predefined variance schedule,
    \item $t = 1, \dots, T$.
\end{itemize}

Define:

\begin{equation}
\alpha_t = 1 - \beta_t,
\quad
\bar{\alpha}_t = \prod_{s=1}^{t} \alpha_s
\end{equation}

Then, sampling directly from timestep $t$ can be written in closed form:

\begin{equation}
q(\mathbf{x}_t \mid \mathbf{x}_0) =
\mathcal{N}\left(
\sqrt{\bar{\alpha}_t}\mathbf{x}_0,
(1 - \bar{\alpha}_t) I
\right)
\end{equation}

Equivalently:

\begin{equation}
\mathbf{x}_t =
\sqrt{\bar{\alpha}_t}\mathbf{x}_0 +
\sqrt{1 - \bar{\alpha}_t}\boldsymbol{\epsilon},
\quad
\boldsymbol{\epsilon} \sim \mathcal{N}(0,I)
\end{equation}

As $t \to T$, $\mathbf{x}_T$ approaches pure Gaussian noise.

\subsection{Reverse Denoising Process}

The goal is to learn a reverse process:

\begin{equation}
p_\theta(\mathbf{x}_{t-1} \mid \mathbf{x}_t)
\end{equation}

parameterized as:

\begin{equation}
p_\theta(\mathbf{x}_{t-1} \mid \mathbf{x}_t)
=
\mathcal{N}\left(
\boldsymbol{\mu}_\theta(\mathbf{x}_t, t),
\Sigma_\theta(\mathbf{x}_t, t)
\right)
\end{equation}

This is modeled using a neural network (typically a U-Net).

\subsection{Variational Objective}

Training maximizes a variational lower bound on log-likelihood:

\begin{equation}
\log p_\theta(\mathbf{x}_0)
\ge
\mathcal{L}_{\text{VLB}}
\end{equation}

The simplified training objective reduces to noise prediction:

\begin{equation}
\mathcal{L}_{\text{simple}} =
\mathbb{E}_{\mathbf{x}_0, \boldsymbol{\epsilon}, t}
\left[
\left\|
\boldsymbol{\epsilon}
-
\boldsymbol{\epsilon}_\theta(\mathbf{x}_t, t)
\right\|_2^2
\right]
\end{equation}

Thus, the model learns to predict the added noise at each timestep.

\subsection{Score Matching Interpretation}

Diffusion models can be interpreted as learning the score function:

\begin{equation}
\nabla_{\mathbf{x}_t}
\log q(\mathbf{x}_t)
\end{equation}

Using denoising score matching:

\begin{equation}
\mathcal{L}_{\text{DSM}} =
\mathbb{E}
\left[
\left\|
s_\theta(\mathbf{x}_t, t)
-
\nabla_{\mathbf{x}_t}
\log q(\mathbf{x}_t \mid \mathbf{x}_0)
\right\|_2^2
\right]
\end{equation}

This connects diffusion models to stochastic differential equations.

\subsection{Sampling Procedure}

Generation starts from Gaussian noise:

\begin{equation}
\mathbf{x}_T \sim \mathcal{N}(0, I)
\end{equation}

Then iteratively sample:

\begin{equation}
\mathbf{x}_{t-1}
=
\frac{1}{\sqrt{\alpha_t}}
\left(
\mathbf{x}_t
-
\frac{\beta_t}{\sqrt{1 - \bar{\alpha}_t}}
\boldsymbol{\epsilon}_\theta(\mathbf{x}_t, t)
\right)
+
\sigma_t \mathbf{z}
\end{equation}

where $\mathbf{z} \sim \mathcal{N}(0,I)$ if $t > 1$.

\subsection{Classifier Guidance}

To condition generation on label $y$, guidance modifies the score:

\begin{equation}
\nabla_{\mathbf{x}_t}
\log p(\mathbf{x}_t \mid y)
=
\nabla_{\mathbf{x}_t}
\log p(\mathbf{x}_t)
+
\nabla_{\mathbf{x}_t}
\log p(y \mid \mathbf{x}_t)
\end{equation}

Classifier-free guidance approximates this without an external classifier.

\subsection{DDIM: Deterministic Sampling}

DDIM modifies sampling to:

\begin{equation}
\mathbf{x}_{t-1}
=
\sqrt{\bar{\alpha}_{t-1}}
\hat{\mathbf{x}}_0
+
\sqrt{1 - \bar{\alpha}_{t-1}}
\boldsymbol{\epsilon}_\theta(\mathbf{x}_t, t)
\end{equation}

allowing faster deterministic sampling.

\subsection{Advantages Over GANs and VAEs}

Diffusion models provide:

\begin{itemize}
    \item Stable training (no adversarial minimax instability)
    \item Full mode coverage (no mode collapse)
    \item Explicit likelihood optimization
    \item High-fidelity image synthesis
\end{itemize}

Compared to VAEs:
\begin{itemize}
    \item Avoid blurry reconstructions
    \item Do not require strong latent Gaussian compression
\end{itemize}

Compared to GANs:
\begin{itemize}
    \item Better diversity
    \item More stable convergence
    \item Improved anatomical consistency
\end{itemize}

\subsection{Relevance to Medical Image Generation}

For chest X-ray generation:

\begin{itemize}
    \item Diffusion models preserve fine-grained anatomical structures.
    \item Enable controlled pathology synthesis.
    \item Provide better diversity for rare disease modeling.
\end{itemize}

This explains their recent dominance in medical image synthesis compared to VAEs and GANs.s